\section{Conclusion}\label{sec:8}\label{conclusion}

Certified comparison can stop early when the evidence separates the finite mean from the relevant practical boundaries. At one-item resolution, many admissible unseen coordinates can each change the answer, forcing substantial evaluation even without observed variation. Between those regimes, disagreement, edit capacity and inference design matter; distance alone does not determine cost. The theory establishes the fine-resolution obstruction, and public-data replays show its practical counterpart: selected boundary-close comparisons can consume almost the whole benchmark. Quantifying the optimal cost throughout the intervening regime remains open.
